\subsubsection{模型的求解}\label{sec:q3-3}

模型里没有一个可以一次解出的优化问题：覆盖核验、区域更新和清除判据是确定的几何计算，选探测点和排路线各是一个小规模的搜索，把它们按规则串起来、每执行一步用真实反馈更新一次，就是求解过程。流程见图\ref{fig:q3-flow}，步骤归纳为算法\ref{alg:q3-search}。整个过程只用四种真实反馈：示向度、$5$ m 内的近场反馈、无信号，以及清除成功与否。几何部分保证不漏源、不裁掉真实源、判定能清除时一定能清除；选点代价、路线搜索和提前光学的面积占比只影响快慢，算错了最多多花时间，不会漏掉源。

\begin{center}
\begin{minipage}{\linewidth}
\centering
\QThreeFlowFigure
\captionof{figure}{问题三的求解流程。每执行一个任务就重新排路线；处理源时先看能否直接清除，再看是否满足提前光学的条件，最后才选点测向。}
\label{fig:q3-flow}
\end{minipage}
\end{center}

\par\medskip
\noindent\begin{minipage}{\linewidth}
\small
\refstepcounter{algorithm}\label{alg:q3-search}
\textbf{算法\thealgorithm：多源搜索与清除的滚动规划}\par\smallskip
\noindent\textbf{输入：}目标圆域 $\Omega$、$20$ 个频道、源数范围 $10$ 到 $16$、$7$ 个检测站、设备动作费用
\begin{enumerate}
\setlength{\itemsep}{2pt}
\setlength{\parskip}{0pt}
\item 从原点出发，$20$ 个频道全部记为未发现，$7$ 个站全部记为未去；
\item\label{step:q3-check} 对每个已知源，若当前位置离其区域所有顶点都不超过 $20$ m，就地清除；若已清除 $16$ 个频道，或未去的站为空且已知源全部清除，结束；
\item 若还有未发现频道，尝试剪掉多余的站、把未去的站往路线上挪，每次改动都重新通过覆盖核验；
\item 把未去的站和已知源（用包围圆圆心代表）排成从当前位置出发的最短路线，取第一个任务，记下第二个任务的位置；
\item 首任务是检测站：前往该站，从当前频道开始把未发现频道逐个测一遍；对满足补测条件的已知源再测一次；把该站记为已扫描；若未发现频道在所有已扫描站都无信号且未去的站为空，判这些频道无源；回到步骤\ref{step:q3-check}；
\item 首任务是源：若包围圆半径不超过 $19.5$ m，走到兼顾下一任务的清除点清除；否则，若该源满足提前光学的资格，走到区域的面积加权重心先试一次清除，失败则原地测向；否则按式\eqref{eq:q3-4} 在候选点中选代价最小的点测向（局部测向已满 $4$ 次或半径小于 $22$ m 时改测圆心）；用真实示向度和无信号记录更新区域；更新后若半径不超过 $19.5$ m 则立即清除；对满足补测条件的其他已知源再测一次；回到步骤\ref{step:q3-check}。
\end{enumerate}
\end{minipage}
\par\medskip

程序用 Python 实现，几何部分只用 NumPy，选点代价和路线搜索用 Numba 编译。在本地模拟器上，一局从进入到退出的计算时间平均 $0.057$ s，远小于机器狗的动作时间，因此每一步都重新规划不构成负担。运行中有三处检查：在圆心测向却收不到信号、更新后的区域为空、判定能清除却清除失败，任何一处出现都说明前面的推理和实际反馈矛盾，程序立即报错停机，不会把矛盾的状态继续用下去；主批 $10000$ 次运行中这三处检查一次也没有触发。
